
% ===
\section{Introduction}
% ===

The relationship between economic activity and the natural environment has long been marked by a persistent tension between productive expansion and ecological preservation. Economic theory does not ignore nature. Land, natural resources, environmental costs, externalities, and common-pool resources have all been extensively discussed in different strands of the literature \cite{smith1776, ricardo1817, pigou2010liberty, coase1960social, hardin1968tragedy, ostrom1990commons, dasgupta2001human}. The problem addressed in this thesis is more specific: in conventional production theory, the environment is rarely treated as an endogenous, dynamic, and strategically relevant component of production decisions. In canonical representations of production, such as the Cobb-Douglas function, output is usually modeled as a function of capital and labor, sometimes with land or natural resources appearing as inputs, but ecological systems are not generally represented as dynamic cofactors that affect future productivity, strategic incentives, and intertemporal welfare \cite{cobb1928, jehle2011}. Instead, the environment often appears as a stock, a constraint, a source of externalities, or an object of corrective regulation.

This theoretical positioning has important implications. When the environment is treated mainly as an external condition of production, environmental preservation tends to be perceived by private agents as a legal obligation, a compliance cost, or a restriction on productive activity. In this setting, the economic agent is frequently placed before a trade-off: to preserve ecological assets in the interest of collective welfare or to expand productive use in the interest of private profitability. The central difficulty is not merely ethical or legal, but economic. If preservation does not alter the agent's expected payoff in a stable and predictable way, compliance will tend to depend on coercion, moral commitment, reputational pressure, or fragmented compensatory mechanisms.

This problem becomes particularly visible in the case of land use. Standard economic models usually define production as a function of land or natural resources, capital, and labor. Although ecological conditions are directly or indirectly embedded in these factors, the environment itself is often treated as exogenous to the productive process \cite{jehle2011}. Ecological functions such as soil fertility, water regulation, biodiversity, carbon sequestration, resilience, and regenerative capacity may affect productivity over time, but they are not always represented as variables that enter the agent's decision-making structure. As a result, environmental degradation is frequently addressed only after it occurs, through taxes, fines, restrictions, compensations, or restoration obligations.

The main institutional responses to environmental degradation reflect this logic. Pigouvian taxation seeks to internalize negative externalities by increasing the private cost of socially harmful activities \cite{pigou2010liberty}. Coasian bargaining, in turn, emphasizes that externalities may be addressed through the definition and exchange of property rights when transaction costs are sufficiently low and parties can negotiate over the allocation of resources \cite{coase1960social}. Hardin's formulation of the tragedy of the commons highlights the risk that individually rational behavior may lead to the overuse of shared resources when access is weakly regulated or when collective constraints are absent \cite{hardin1968tragedy}. Common-pool resource theory, especially in Ostrom's work, shows that, under certain institutional conditions, communities may self-govern shared resources without necessarily relying only on privatization or centralized state control \cite{ostrom1990commons, ostrom2002commonpool}. Command-and-control regulation imposes restrictions, prohibitions, and sanctions on environmentally harmful conduct, while carbon markets, environmental certifications, payments for environmental services, and fiscal incentives attempt to create positive economic rewards for conservation or reduced degradation.

These mechanisms have contributed significantly to environmental policy design. However, they also present limitations. Negative incentives, such as fines, embargoes, and legal restrictions, may deter illegal behavior, but they do not necessarily make preservation economically attractive. Positive incentives, such as carbon credits, certifications, subsidies, and payments for environmental services, may generate additional revenue, but they often remain fragmented, uncertain, or detached from the ecological unit affected by production. In many cases, environmental preservation remains an external condition attached to production rather than an endogenous component of the production process itself.

A second limitation concerns the spatial design of environmental compensation. Ecological degradation is territorially specific. A loss of native vegetation, biodiversity, hydrological function, soil resilience, or carbon storage occurs in a concrete ecological unit. Yet many compensation mechanisms operate through standardized, centralized, or globally tradable instruments. A producer may offset environmental impacts in a location different from the territory affected by production. This may satisfy an accounting or regulatory requirement, but it does not necessarily restore the ecological functions of the affected biome, watershed, ecosystem, or productive region. The problem, therefore, is not only whether compensation occurs, but where it occurs and whether it preserves the ecological conditions that sustain local production and long-term resilience.

A third limitation concerns time. Environmental degradation and preservation are intertemporal processes. The benefits of exploiting natural resources may be immediate, private, and measurable, whereas the benefits of preservation are often delayed, uncertain, collective, and spatially diffuse. This temporal asymmetry weakens the private incentive to preserve, especially when the agent faces credit constraints, competitive pressure, incomplete information, weak enforcement, or political uncertainty. As a result, even agents who understand the long-term importance of ecological preservation may still choose short-term expansion if the expected private payoff of degradation exceeds the expected private payoff of preservation.

These limitations suggest that the opposition between production and preservation may be less a natural incompatibility than a consequence of theoretical and institutional misalignment. If ecological variables remain external to production theory, and if environmental policy operates mainly as a corrective layer imposed after private decisions have already been made, preservation will tend to appear as a burden. Conversely, if ecological preservation is incorporated into the productive calculus of agents, it may become possible to design institutional arrangements in which preserving environmental assets becomes compatible with productivity, profitability, and long-term welfare.

This thesis develops this possibility through the provisional concept of \emph{economic commensalism}. The term is not introduced here as a fully established theory, nor as a direct transposition of a biological category into economics. Rather, it is used as an analytical hypothesis to investigate whether there can exist institutional and productive arrangements in which economic agents benefit from production while ecological systems are maintained within functional and regenerative thresholds. In this sense, commensalism is not understood as the absence of environmental impact, but as a configuration in which productive activity does not systematically undermine the ecological conditions that sustain future production.

The concept is therefore presented as a research object to be developed, formalized, and empirically assessed throughout the thesis. At this stage, \emph{economic commensalism} names a theoretical possibility: the possibility that environmental preservation may become part of the agent's rational decision-making process when ecological variables are endogenized into production theory and when incentives are designed according to spatial, intertemporal, and strategic constraints. The concept does not deny the relevance of regulation, enforcement, markets, or collective governance. Rather, it seeks to reorganize their role around a central question: under what conditions can environmental preservation become individually rational, strategically stable, and socially desirable?

The proposed framework differs from predatory and extractivist patterns of resource use. Predation refers to the destructive, illegal, or unsustainable depletion of natural resources, often disregarding property rights, ecological limits, and collective welfare. Extraction, in a neutral sense, refers to the use or removal of natural resources as part of productive activity. Extractivism, by contrast, is understood here as a productive and institutional regime in which resource use becomes structurally disconnected from ecological regeneration, future productivity, and the alignment between private and social payoffs. The distinction is important because not every productive use of nature is necessarily predatory; however, productive use becomes extractivist when it systematically undermines the ecological conditions upon which future production depends.

From this perspective, the central issue is not simply whether natural resources should be used or preserved. The relevant question is how institutions, incentives, and production theory can distinguish between productive use that maintains ecological functionality and productive use that erodes the basis of future production. Economic commensalism is proposed as a framework for investigating this distinction. It assumes, as a hypothesis to be examined, that preservation may become economically rational when the benefits of maintaining ecological integrity are linked to the agent's expected returns, access to markets, credit conditions, production rights, fiscal treatment, reputational gains, technological monitoring, and long-term productivity.

The Brazilian state of Mato Grosso provides a particularly relevant empirical and conceptual context for this investigation. The state is one of Brazil's main agricultural frontiers and contains portions of three major biomes: the Amazon, the Cerrado, and the Pantanal. It is therefore located at the intersection of agricultural expansion, deforestation pressure, biodiversity conservation, carbon-related opportunities, and institutional disputes over land use. The ecological heterogeneity of Mato Grosso makes it especially useful for analyzing whether environmental governance should be designed according to broad political boundaries or according to ecological units such as biomes, watersheds, ecoregions, and transition zones.

Brazilian environmental law already reflects, at least partially, the importance of ecological differentiation. The Brazilian Forest Code establishes different legal requirements for native vegetation preservation depending on the region and biome in which rural properties are located \cite{brasil2012florestal}. However, legal restrictions alone do not solve the incentive problem. If producers bear immediate opportunity costs from preservation while receiving uncertain, delayed, or diffuse benefits, the law may generate resistance, circumvention, or dependence on costly enforcement. This reinforces the need to investigate institutional mechanisms capable of making preservation compatible with the agent's economic calculus.

The thesis therefore rests on four central premises. First, the environment can be represented as an endogenous cofactor of production, rather than merely as an externality or constraint. Second, many existing environmental instruments fail to alter the structural private payoff that drives degradation. Third, preservation becomes stable only when it is compatible with individual rationality and strategic interaction among agents. Fourth, such compatibility requires governance mechanisms that are locally anchored in ecological units while remaining connected to broader systems of verification, enforcement, and institutional credibility.

The general research problem can be stated as follows: can environmental preservation be theoretically reframed not merely as a constraint on production, but as an endogenous and incentive-compatible component of economic decision-making? More specifically, the thesis asks whether private productivity and profitability can be reconciled with socially desirable environmental preservation through institutional arrangements that are spatially grounded, intertemporally coherent, and strategically stable.

The main hypothesis is that, when ecological variables are incorporated into production theory and when incentives are designed to reflect spatial, intertemporal, and behavioral constraints, environmental preservation can become compatible with individual rationality rather than dependent solely on coercion, moral commitment, or fragmented compensation. In this sense, the thesis investigates whether economic commensalism can provide a theoretical basis for aligning private incentives with ecological preservation without necessarily reducing the long-term productive viability of economic agents.

The objective of the thesis is to develop the conceptual, formal, and empirical foundations of this framework. The theoretical objective is to clarify the limitations of approaches that treat environmental degradation mainly as an externality, a property-rights problem, or a regulatory failure. The formal objective is to model the interaction between production, preservation, uncertainty, bounded rationality, local ecological feedback, and strategic equilibrium. The empirical objective is to evaluate the feasibility of the proposed framework using Mato Grosso as a case in which agricultural productivity, deforestation, ecological heterogeneity, and environmental governance interact in a concrete institutional setting.

The thesis is organized into three main chapters. Chapter 1 develops the theoretical foundations of economic commensalism by discussing the relationship between production theory, environmental preservation, extractivism, predation, common-pool resources, and incentive-compatible governance. It also introduces a preliminary formal representation of the environment as an endogenous cofactor of production, in order to clarify the conceptual mechanism through which ecological preservation may affect productivity, incentives, and intertemporal welfare. Chapter 2 extends this preliminary representation into a dynamic mathematical model that incorporates environmental variables into production decisions, with attention to intertemporal choice, uncertainty, learning, ecological feedback, and strategic interaction. Chapter 3 provides an empirical application to Mato Grosso, examining whether the theoretical and formal mechanisms proposed in the previous chapters are consistent with observed economic and environmental dynamics.

By reframing environmental preservation as a potentially endogenous component of economic rationality, this thesis seeks to contribute to environmental and resource economics, institutional economics, and public policy design. Its central contribution is not to reject existing instruments such as regulation, enforcement, carbon markets, payments for environmental services, or community governance, but to investigate how these instruments can be reorganized around a more precise theoretical problem: the need to transform preservation from an external burden into a strategically rational and locally grounded component of production.

\newpage

% fim %
